\chapter{Organiser des données}\label{ch:g7:stats}

Avant de calculer quoi que ce soit, les données doivent être
recueillies, comptées et présentées. Ce chapitre concerne la
transformation d'une liste désordonnée d'observations en un tableau ou
une image claire~: \omterm{def:g7:stats:frequency}{effectifs}, \omterm{def:g7:stats:frequency}{fréquences}, données regroupées,
diagrammes en barres et diagrammes circulaires. Résumer des données
par un seul nombre (moyenne, médiane) est le sujet du
\cref{ch:g9:statproba}.

\section{Effectifs et fréquences}

\begin{definition}[Effectif et fréquence]\label{def:g7:stats:frequency}
Pour chaque valeur possible d'une enquête, son
\emph{effectif}\index{effectif} est le nombre de fois où elle
apparaît~; sa \emph{fréquence}\index{fréquence} est l'effectif divisé
par le nombre total d'observations~:
\[
\text{fréquence} = \frac{\text{effectif}}{\text{total}} .
\]
Les fréquences peuvent s'écrire en \omterm{def:g6:fractions:def}{fractions}, en décimaux ou en
pourcentages, et elles s'additionnent toujours à $1$ (c'est-à-dire
$100\,\%$).
\end{definition}

\begin{example}\label{ex:g7:stats:frequency}
Sport préféré des $25$ élèves d'une classe, de la liste brute au
tableau~:
\begin{center}
\renewcommand{\arraystretch}{1.3}
\begin{tabular}{l|cccc|c}
sport & football & natation & tennis & judo & total \\
\hline
\omterm{def:g7:stats:frequency}{effectif} & $10$ & $6$ & $5$ & $4$ & $25$ \\
\omterm{def:g7:stats:frequency}{fréquence} & $0.40$ & $0.24$ & $0.20$ & $0.16$ & $1$ \\
\end{tabular}
\end{center}
Pour le football~: $\frac{10}{25} = \frac{40}{100} = 0.40 = 40\,\%$.
La ligne des \omterm{def:g7:stats:frequency}{fréquences} doit sommer à $1$ --- un contrôle permanent
et gratuit.
\end{example}

\begin{remark}[Pourquoi des fréquences~?]
Les \omterm{def:g7:stats:frequency}{effectifs} ne se comparent pas d'un groupe à l'autre lorsque les
tailles diffèrent~: $10$ amateurs de football dans une classe de $25$
(soit $40\,\%$) montrent plus d'enthousiasme que $12$ dans une année
scolaire de $120$ ($10\,\%$). Les \omterm{def:g7:stats:frequency}{fréquences} mettent tout à la même
échelle.
\end{remark}

\section{Regrouper des données en classes}

\begin{example}[Classes]\label{ex:g7:stats:classes}
Les tailles (en cm) de $20$ élèves s'étendent de $148$ à $172$.
Lister chaque taille séparément donnerait un tableau avec presque
autant de colonnes que d'élèves~; à la place, regrouper en
\emph{classes}~:
\begin{center}
\renewcommand{\arraystretch}{1.3}
\begin{tabular}{l|ccc|c}
taille (cm) & $\intco{145}{155}$ & $\intco{155}{165}$ &
$\intco{165}{175}$ & total \\
\hline
\omterm{def:g7:stats:frequency}{effectif} & $5$ & $9$ & $6$ & $20$ \\
\omterm{def:g7:stats:frequency}{fréquence} & $0.25$ & $0.45$ & $0.30$ & $1$ \\
\end{tabular}
\end{center}
Chaque taille est comptée dans exactement une classe ($155$ va dans
$\intco{155}{165}$, pas dans les deux --- c'est pourquoi les
intervalles sont semi-ouverts). Le regroupement perd du détail mais
gagne en lisibilité.
\end{example}

\section{Diagrammes}

\begin{method}[Choisir et tracer un diagramme]\label{met:g7:stats:charts}
\begin{enumerate}
  \item \emph{Diagramme en barres}~: une barre par valeur ou classe,
        hauteur proportionnelle à l'\omterm{def:g7:stats:frequency}{effectif} (ou à la \omterm{def:g7:stats:frequency}{fréquence}). Bon
        pour comparer des valeurs.
  \item \emph{Diagramme circulaire}~: un secteur par valeur, angle
        proportionnel à la \omterm{def:g7:stats:frequency}{fréquence} --- le cercle entier
        ($360^\circ$) représente le total, donc
        \[
        \text{angle} = \text{fréquence} \times 360^\circ .
        \]
        Bon pour montrer les parts d'un tout.
  \item Dans les deux cas~: titre, légendes, et une graduation ou les
        pourcentages écrits sur le diagramme.
\end{enumerate}
\end{method}

\begin{example}\label{ex:g7:stats:pie}
Angles pour les sports de \cref{ex:g7:stats:frequency}~: football
$0.40 \times 360 = 144^\circ$~; natation $0.24 \times 360 =
86.4^\circ$~; tennis $72^\circ$~; judo $57.6^\circ$. Vérification~:
$144 + 86.4 + 72 + 57.6 = 360$.
\end{example}

\begin{omfigure}
\begin{tikzpicture}
\begin{axis}[
  width=6.6cm, height=5.2cm,
  ybar, bar width=14pt,
  symbolic x coords={football,natation,tennis,judo},
  xtick=data,
  x tick label style={font=\footnotesize, rotate=25, anchor=east},
  ymin=0, ymax=12,
  ytick={2,4,6,8,10},
  ylabel={effectif},
  label style={font=\small},
  tick label style={font=\small}]
  \addplot[fill=omDef!30, draw=omDef] coordinates
    {(football,10) (natation,6) (tennis,5) (judo,4)};
\end{axis}
\end{tikzpicture}\hspace{0.8cm}%
\begin{tikzpicture}[scale=1.05]
  \fill[omDef!45] (0,0) -- (0:1.5) arc (0:144:1.5) -- cycle;
  \fill[omThm!35] (0,0) -- (144:1.5) arc (144:230.4:1.5) -- cycle;
  \fill[omMeth!35] (0,0) -- (230.4:1.5) arc (230.4:302.4:1.5) -- cycle;
  \fill[omProp!35] (0,0) -- (302.4:1.5) arc (302.4:360:1.5) -- cycle;
  \draw (0,0) circle (1.5);
  \foreach \a in {0,144,230.4,302.4} \draw (0,0) -- (\a:1.5);
  \node[font=\footnotesize] at (72:0.95) {foot.};
  \node[font=\footnotesize] at (187:0.95) {nat.};
  \node[font=\footnotesize] at (266:0.95) {tennis};
  \node[font=\footnotesize] at (331:0.95) {judo};
\end{tikzpicture}

{\small Les mêmes données en diagramme en barres (comparer les
hauteurs) et en diagramme circulaire (comparer les parts de la
classe). Le football occupe $144^\circ$, soit $40\,\%$ du cercle.}
\end{omfigure}

\begin{example}[Lire avec un œil critique]\label{ex:g7:stats:critical}
Un diagramme dont l'axe vertical commence à $9$ au lieu de $0$ fait
paraître une barre de $10$ dix fois plus haute qu'une barre de $9.1$
--- trompeur visuellement, même si les nombres sont honnêtes. Premier
réflexe face à tout diagramme~: vérifier où l'axe commence et ce que
vaut une graduation (\cref{met:g6:propdata:read}).
\end{example}

\section{Exercices}

\begin{exercise}[$\star$]\label{exo:g7:stats:1}
Voici les notes de $20$ élèves~: $12$, $15$, $9$, $12$, $14$, $9$,
$12$, $15$, $11$, $12$, $9$, $14$, $12$, $11$, $15$, $9$, $12$, $14$,
$11$, $12$. Construire le tableau des \omterm{def:g7:stats:frequency}{effectifs} (une colonne par
note).
\end{exercise}

\begin{exercise}[$\star$]\label{exo:g7:stats:2}
Compléter le tableau de \cref{exo:g7:stats:1} d'une ligne de
\omterm{def:g7:stats:frequency}{fréquences} (en \omterm{def:g6:fractions:def}{fractions} de $20$, puis en pourcentages). Vérifier la
\omterm{def:g1:addition:def}{somme}.
\end{exercise}

\begin{exercise}[$\star$]\label{exo:g7:stats:3}
Dans une enquête auprès de $50$ familles, $30$ ont une voiture.
Quelle est la \omterm{def:g7:stats:frequency}{fréquence} de «~une voiture~», en décimal et en
pourcentage~? Dans une autre enquête, $45$ familles sur $90$ ont une
voiture~: laquelle montre la proportion la plus élevée~?
\end{exercise}

\begin{exercise}[$\star$]\label{exo:g7:stats:4}
Les \omterm{def:g2:measure:units}{masses} (kg) de $15$ chiens~: $8$, $23$, $31$, $12$, $9$, $27$,
$35$, $14$, $22$, $18$, $29$, $11$, $25$, $33$, $16$. Les regrouper
dans les classes $\intco{5}{15}$, $\intco{15}{25}$, $\intco{25}{35}$,
$\intco{35}{45}$ et donner les \omterm{def:g7:stats:frequency}{effectifs}.
\end{exercise}

\begin{exercise}[$\star$]\label{exo:g7:stats:5}
Tracer le diagramme en barres du tableau~:
\begin{center}
\renewcommand{\arraystretch}{1.2}
\begin{tabular}{l|cccc}
animaux & $0$ & $1$ & $2$ & $3$ \\
\hline
familles & $8$ & $12$ & $6$ & $4$ \\
\end{tabular}
\end{center}
\end{exercise}

\begin{exercise}[$\star$]\label{exo:g7:stats:6}
Pour le tableau de \cref{exo:g7:stats:5}, calculer les \omterm{def:g7:stats:frequency}{fréquences} et
les angles du diagramme circulaire de chaque valeur (vérification~:
total $360^\circ$).
\end{exercise}

\begin{exercise}[$\star$]\label{exo:g7:stats:7}
Un diagramme circulaire sur les saisons préférées montre~: été
$162^\circ$, printemps $90^\circ$, automne $54^\circ$, hiver
$54^\circ$. Quel pourcentage a choisi chaque saison~? Si $60$
personnes ont répondu, combien ont choisi l'été~?
\end{exercise}

\begin{exercise}[$\star\star$]\label{exo:g7:stats:8}
Dans la classe A, $12$ élèves sur $30$ viennent à pied~; dans la
classe B, $14$ sur $40$. Quelle classe a la proportion de piétons la
plus élevée~? Justifier avec des \omterm{def:g7:stats:frequency}{fréquences}, pas des \omterm{def:g7:stats:frequency}{effectifs}.
\end{exercise}

\begin{exercise}[$\star\star$]\label{exo:g7:stats:9}
Un magazine imprime un diagramme en barres des ventes mensuelles~:
$9\,800$, puis $10\,000$, puis $10\,100$, avec l'axe vertical
commençant à $9\,700$. Décrire ce que le lecteur voit, et ce qu'un
axe honnête partant de $0$ montrerait à la place.
\end{exercise}

\begin{exercise}[$\star\star$]\label{exo:g7:stats:10}
Un tableau de \omterm{def:g7:stats:frequency}{fréquences} a trois valeurs de \omterm{def:g7:stats:frequency}{fréquences} $0.35$, $0.4$
et $f$. Trouver $f$. Si l'\omterm{def:g7:stats:frequency}{effectif} total est $80$, donner les trois
\omterm{def:g7:stats:frequency}{effectifs}.
\end{exercise}

\begin{exercise}[$\star\star\star$]\label{exo:g7:stats:11}
Dans un établissement, $55\,\%$ des élèves sont des filles. Parmi les
filles, $40\,\%$ mangent à la cantine~; parmi les garçons, $60\,\%$
le font. Sur $400$ élèves, combien mangent à la cantine~? (Calculer
les quatre tailles de groupes étape par étape.) Quel pourcentage
global cela représente-t-il~?
\end{exercise}

\section{Problème~: comment les données mentent --- et comment les
prendre en flagrant délit}

\begin{problem}[{Devoir maison --- effectifs contre fréquences, le
sondage qui a trompé tout un pays, et le paradoxe des deux hôpitaux}]
\label{pb:g7:stats:1}
Les nombres ne mentent pas, mais on peut s'en servir pour induire
en erreur~: un \omterm{def:g7:stats:frequency}{effectif} énorme peut cacher une petite \omterm{def:g7:stats:frequency}{fréquence}, une
enquête gigantesque peut ne rien valoir, et --- le plus étrange ---
un hôpital peut battre son rival sur \emph{toutes} les catégories de
patients tout en perdant sur les chiffres d'ensemble. Ce problème
déclenche les trois pièges, armé des seuls \omterm{def:g7:stats:frequency}{effectifs} et \omterm{def:g7:stats:frequency}{fréquences} de ce
chapitre (\cref{def:g7:stats:frequency}).

\medskip
\noindent\textbf{Partie I --- Les \omterm{def:g7:stats:frequency}{effectifs} ne sont pas des
\omterm{def:g7:stats:frequency}{fréquences}.}
\begin{enumerate}
  \item L'école A recycle $120$ de ses $400$ briques de jus de
        fruits~; l'école B, $90$ sur $250$. Quelle école a le plus
        grand \emph{\omterm{def:g7:stats:frequency}{effectif}} de briques recyclées~? La plus grande
        \emph{\omterm{def:g7:stats:frequency}{fréquence}}~? Quelle école mérite le prix du
        recyclage~?
  \item Un nouvel emploi du temps fait l'objet d'une enquête~: $18$
        des $30$ professeurs l'apprécient, et $210$ des $600$ élèves
        aussi. Calculer les deux \omterm{def:g7:stats:frequency}{fréquences}, puis la \omterm{def:g7:stats:frequency}{fréquence} sur
        l'ensemble des $630$ personnes. Pourquoi le chiffre global
        est-il si proche de celui des élèves~?
  \item Un diagramme circulaire publié montre des secteurs étiquetés
        $45\,\%$, $30\,\%$, $20\,\%$ et $10\,\%$. Sans aucune autre
        information, comment sait-on qu'une erreur a été commise~?
  \item Un tableau de \omterm{def:g7:stats:frequency}{fréquences} portant sur $40$ observations est à
        \omterm{def:g3:division:half}{moitié} effacé~: \omterm{def:g7:stats:frequency}{effectifs} $14$, $10$, $?$, $?$ et \omterm{def:g7:stats:frequency}{fréquences}
        $0.35$, $0.25$, $0.15$, $?$. Reconstituer les cases
        manquantes.
  \item Une enquête aboutit aux \omterm{def:g7:stats:frequency}{fréquences} $0.40$, $0.35$ et $0.25$.
        Calculer les trois angles du diagramme circulaire
        (\cref{met:g7:stats:charts}) et vérifier qu'ils referment le
        cercle.
\end{enumerate}

\medskip
\noindent\textbf{Partie II --- Le sondage qui a trompé tout un pays.}
En 1936, un magazine américain envoya dix millions de bulletins à des
adresses tirées des annuaires téléphoniques et des listes
d'immatriculation d'automobiles, reçut $2.4$ millions de réponses, et
prédit une victoire écrasante du candidat Landon. Un jeune
statisticien, George Gallup, n'interrogea qu'une cinquantaine de
milliers de personnes --- choisies pour ressembler à la population
entière --- et prédit l'inverse. Roosevelt l'emporta très largement.
\begin{enumerate}[resume]
  \item Quelle \omterm{def:g7:stats:frequency}{fréquence} des bulletins envoyés est revenue~?
  \item En 1936, le téléphone et l'automobile étaient des \omterm{def:g2:mult:def}{produits} de
        luxe. Expliquer en une ou deux phrases pourquoi
        l'échantillon du magazine, pourtant énorme, était condamné
        --- et quelle leçon de la question 1 il répète à l'échelle
        d'un pays.
  \item Un modèle réduit. Une ville compte $1\,000$ électeurs aisés,
        dont $70\,\%$ soutiennent L, et $9\,000$ électeurs modestes,
        dont $30\,\%$ soutiennent L. Calculer le soutien réel de L
        dans la ville. Un «~sondage par annuaire~» touche $500$
        électeurs aisés et $500$ modestes~: quel soutien
        prédit-il~?
  \item Le piège silencieux~: une entreprise a $200$ clients
        mécontents et $800$ clients satisfaits. Elle les interroge
        tous~; $40\,\%$ des mécontents répondent, mais seulement
        $10\,\%$ des satisfaits. Calculer le nombre de réponses de
        chaque sorte, et la \omterm{def:g7:stats:frequency}{fréquence} du mécontentement \emph{parmi
        les réponses}. Comparer avec la \omterm{def:g7:stats:frequency}{fréquence} réelle.
  \item Réparer le sondage de la question 8~: garder $1\,000$
        entretiens, mais les répartir dans les proportions réelles
        de la ville. Que prédit le sondage réparé~?
\end{enumerate}

\medskip
\noindent\textbf{Partie III --- Le paradoxe des deux hôpitaux.}
Deux hôpitaux publient leurs \omterm{def:g7:stats:frequency}{effectifs} de guérisons, répartis par
gravité~:
\begin{center}
\renewcommand{\arraystretch}{1.3}
\begin{tabular}{l|cc}
 & cas bénins & cas graves \\
\hline
hôpital A & $90$ guéris sur $100$ & $280$ guéris sur $400$ \\
hôpital B & $340$ guéris sur $400$ & $65$ guéris sur $100$ \\
\end{tabular}
\end{center}
\begin{enumerate}[resume]
  \item Calculer la \omterm{def:g7:stats:frequency}{fréquence} de guérison de l'hôpital A pour les cas
        bénins, pour les cas graves, et globalement (sur les $500$
        patients).
  \item Faire les trois mêmes calculs pour l'hôpital B. Quel hôpital
        l'emporte sur les cas bénins~? Sur les cas graves~?
  \item Comparer les deux \omterm{def:g7:stats:frequency}{fréquences} globales. Énoncer clairement la
        chose étrange qui vient de se produire.
  \item Expliquer le tour de passe-passe~: comparer les
        \emph{compositions} de patientèle des deux hôpitaux, et dire
        pourquoi la \omterm{def:g7:stats:frequency}{fréquence} globale peut trahir les deux \omterm{def:g7:stats:frequency}{fréquences}
        par catégorie. Avec un cas grave, quel hôpital
        faudrait-il choisir~?
  \item Un titre de journal annonce~: «~L'hôpital B a le meilleur
        taux de guérison~: $81\,\%$ contre $74\,\%$.~» Rédiger la
        lettre de deux phrases au rédacteur en chef que tout ce
        problème a appris à écrire --- une phrase sur ce que le titre
        ignore, une phrase sur la morale générale (comparer ce qui
        est comparable, et se méfier des poids).
\end{enumerate}
\end{problem}
